\section{Motivation}\label{sec:motivation}

% deeper into the arguments

\subsection{What's Missing: ...}\label{subsec:what-is-missing}

% Talk about the gap, talk about why CPU has its limit
% A Logical Operator?
% An end to end system?

TBD

\subsection{The Transition: Why Now}\label{subsec:the-transition}

% The shift we observe, and the questions that each quadrant answers

\begin{figure}
  \centering
  \begin{tikzpicture}[
      font=\small\sffamily,
      cell/.style={draw, rounded corners=2pt, minimum width=3.4cm,
                   minimum height=1.6cm, align=center},
      open/.style={cell, dashed, line width=0.9pt, fill=black!8},
      ax/.style={font=\small\sffamily\bfseries},
      flow/.style={-{Latex[length=2.2mm]}, thick, black!65},
    ]

    % four quadrants
    \node[cell] (a) at (0,0)      {70s - 80s};
    \node[cell] (b) at (3.8,0)    {2000s};
    \node[cell] (c) at (0,-2.2)   {now};
    \node[open] (d) at (3.8,-2.2) {vector join};

    % arrows
    \draw[flow] (a) -- (b);
    \draw[flow] (c) -- (d);
    \draw[flow] (a) -- (c);

    % column axis (x)
    \node[ax, anchor=south] at ([yshift=3pt]a.north) {OLTP};
    \node[ax, anchor=south] at ([yshift=3pt]b.north) {OLAP};

    % row axis (y)
    \node[ax, rotate=90, anchor=south] at ([xshift=-3pt]a.west) {relational};
    \node[ax, rotate=90, anchor=south] at ([xshift=-3pt]c.west) {vector};
  \end{tikzpicture}
  \caption{Relational data moved from transactional processing (OLTP,
    70s--80s) to analytical processing (OLAP, 2000s). Vector data is at
    the same inflection point: vector search answers one query at a time
    (OLTP, now), while the analytical quadrant, the vector join, is
    largely unexplored.}
  \label{fig:quadrant}
\end{figure}


The analytical (OLAP) market emerged and grew into its own class of systems because the analytical workload turned out to be a different problem from transaction processing~\cite{StonebrakerPavloAround24}.
OLTP answers ``what matches this'': find this person's record, this order, this account, one lookup and an answer in milliseconds.
OLAP answers ``what matches what'': what were the total sales by region last quarter, the analytical question that relates and combines tables and summarizes across them.
The analytical half is what turned data into meaningful decision in business nowadays for the past decades, as not only did it grow into its own class of systems but also its own field of study.
This is marked as a shift from answering point queries to analyzing the whole dataset.

We believe vector data sits at the same inflection point (Figure~\ref{fig:quadrant}).
Vector search answers the OLTP question, ``what is similar to this'': given one query vector, return its nearest neighbors.
The OLAP question it answers is ``what is similar to what'': score a whole table of vectors against another.
In other words, a vector join, a many-to-many instead of a one-to-many, and it is the half that is missing.
Few systems or vendors offer vector join operations today, and we believe there is going to be a similar transition of the focus, because the workloads for vector join already exist and are in great demand.

\subsection{The Workloads and Related Work}\label{subsec:the-work}

Vector join is the workload that already appears throughout analytical pipelines under different names.
In machine learning pipelines, semantic deduplication joins a training corpus against itself to find and remove near-duplicate data points; pruning this redundancy has been shown to improve both data efficiency and downstream model quality~\cite{AbbasSemDeDup23}.
In data integration, entity resolution matches records across two sources by the similarity of their learned embeddings rather than by exact keys, which turns a classically hard matching problem into an embedding similarity join~\cite{SantanaEMJoin25, SantanaRibeiro23}.
In retrieval and recommendation, an entire batch of user or query embeddings is scored against a catalog or document corpus at once for offline candidate generation, evaluation, or the mining of hard negatives that train the next embedding model~\cite{HuangFBRetrieval20, XiongANCE21, ThakurBEIR21}.
In every case the work is inherently batched.

The demand is now surfacing in native SQL language itself.
Apache Spark has a standing proposal for a \texttt{NEAREST BY} top-$k$ ranking join~\cite{ApacheSparkJIRA56395}, Databricks SQL already ships a \texttt{NEAREST BY} clause~\cite{DatabricksNearestBy26}, and recently DuckDB~\cite{DuckDBNearestByPR24137}.
The declarative surface for vector joins is beginning to standardize, which sharpens the open question: how to \emph{execute} one efficiently.

\subsection{A Vector Join Is a Relational Join}\label{subsec:why-database}

Stonebraker and Pavlo argue that vector search is a feature, not an architecture, and that relational engines will simply absorb it~\cite{StonebrakerPavloAround24}.
The vector join is where that absorption pays off, because it is a relational join, for the following reason.
Its join predicate is a non-equi condition that ranks candidate pairs by distance (unlike a hash join, which is a equi-join by matching index), and it is followed by a reduction that selects which ranked pairs survive.
These are the vector analogues of the reductions that OLAP engines already optimize, and the \texttt{NEAREST BY} proposals encode exactly this structure~\cite{ApacheSparkJIRA56395, DatabricksNearestBy26}.

Treating it as a join brings the standard toolbox to work with.
Reductions can be pushed into the join so that only surviving pairs are ever materialized, in the same spirit as pushing a limit or an aggregate below a join.
More importantly, the join can be \emph{partitioned}.
A hash join partitions its inputs by a hash of the join key so that only co-partitioned rows are compared; a vector join partitions them by \emph{semantic} proximity, clustering both sides so that a query row is compared only against corpus rows in nearby clusters.
This is the IVF-style approximate search recast as a partitioned join, and it is what makes the approximate option available using the same relational machinery.
Whether to use that option is a planning decision.
The planner reads the cardinality of the intermediate result reaching the join and picks accordingly: an exact brute-force join at high selectivity, an approximate join at low selectivity.

\subsection{Why the GPU}\label{subsec:why-gpu}

% This section we can show more of our understanding about the type of workloads and use cases we've identified and show that GPU is the perfect fit for this kind of even-more-compute-intensive workload.

Relational OLAP is the heavier sibling of OLTP, and the intuition carries an expectation that the analytical vector workload should be heavier too, but that is what makes it a better fit for the GPU\@.
The core of a vector join is the computation of pairwise distances between a batch of query vectors and a batch of corpus vectors, which for the common metrics reduces to a dense matrix product.
A single-query search leaves the GPU memory-bound and dominated by launch overhead; batching the query side turns it into a compute-bound GEMM that saturates the tensor cores.
In other words, we achieve higher arithmetic intensity from having batched inputs; it does more work but runs far more efficiently per unit of work.
The core of a vector join reduces to a dense matrix product.
One query at a time barely occupies the GPU; scoring a whole table of queries at once turns that into a single large matrix multiple, which is exactly what the hardware is built to run fast.

Two things follow from this mapping, one expected and one not.
The expected payoff is throughput.
A single GPU delivers what a join needs by reusing the same on-device search primitives that libraries such as cuVS~\cite{NvidiaCuVS26} already provide.
The surprising one is exactness.
A tiled brute-force join computes the distance matrix in blocks and never materializes it in full, so the answer stays exact where CPU systems had to settle for approximate search on cost alone.
Building the operator inside a GPU-native SQL engine~\cite{Yogatama26} keeps it on the device alongside the scans, filters, and aggregations of the surrounding query, so the vector join runs as one stage of the same GPU pipeline.

Moreover, Sirius ranks first on ClickBench~\cite{SiriusClickBench25} on analytical workloads, and it does so on the same join, filter, and aggregate machinery that a vector join reuses.
The distance computation is embarrassingly parallel, since every candidate pair is independent, and the reduction that follows is a standard aggregate.
The vector join therefore lands on the engine's existing strengths naturally.

\subsection{A Motivating Query}\label{subsec:motivating-query}

TBD.
